\chapter{Reporte y Resolución de No Cumplimientos}
\label{sec:reporte-resolucion}

\section{Proceso de Reporte de Defectos}
\label{sec:reporte-defectos}

Los defectos encontrados durante las pruebas se registraron en Jira siguiendo el siguiente flujo:

\begin{enumerate}
    \item \textbf{Reporte:} El tester crea un issue de tipo ``Bug'' en Jira.
    \item \textbf{Clasificación:} El QA Lead asigna severidad (\texttt{Critical}, \texttt{Major}, \texttt{Minor}) y prioridad.
    \item \textbf{Asignación:} El Project Manager asigna el bug al desarrollador responsable.
    \item \textbf{Corrección:} El desarrollador implementa la corrección y mueve el issue a ``In Testing''.
    \item \textbf{Verificación:} El tester re-ejecuta el caso de prueba y verifica la corrección.
    \item \textbf{Cierre:} Si el defecto está corregido, se cierra el issue. En caso contrario, se reabre.
\end{enumerate}

\section{Defectos Encontrados}
\label{sec:defectos-encontrados}

Los defectos encontrados se detallan en el Apéndice~C. A continuación se presenta un resumen:

\begin{longtable}{lccc}
\caption{Resumen de defectos}
\label{tab:resumen-defectos}\\
\toprule
\textbf{ID} & \textbf{Severidad} & \textbf{Estado} & \textbf{Caso Asociado}\\
\midrule
\endfirsthead
\toprule
\textbf{ID} & \textbf{Severidad} & \textbf{Estado} & \textbf{Caso Asociado}\\
\midrule
\endhead
\bottomrule
\endfoot

BUG-002 & Minor & Abierto & CP-04\\

BUG-003 & Critical & Cerrado & CP-14\\

BUG-004 & Major & Cerrado & CP-08\\

BUG-005 & Major & Abierto & CP-08\\

BUG-006 & Minor & Abierto & CP-01\\

BUG-007 & Minor & Cerrado & CP-04\\

BUG-008 & Minor & Abierto & CP-02\\

BUG-009 & Minor & Abierto & CP-06\\

\end{longtable}

\section{Resumen de Resultados}
\label{sec:resumen-resultados}

\begin{itemize}
    \item Total defectos: 8
    \item Defectos cerrados: 3 (BUG-003, BUG-004, BUG-007)
    \item Defectos abiertos: 5 (BUG-002, BUG-005, BUG-006, BUG-008, BUG-009)
    \item Tasa de cierre: 37.5\%
\end{itemize}
